% GENERATED FILE. Edit canonical lessons, then run npm run generate:books.
% canonical-source-hash: fnv1a64:a527e1544790bbbc
% canonical-lessons: BN-C311-nagordola, BN-C311-putulnach, BN-C311-belun, BN-C311-ghora, BN-C311-thasa

\chapter{At the Village Fair}
\label{ch:bn-at-the-village-fair}
\hlchaptermodality{311}

\begin{chapteropening}
\textbf{By the end of this chapter:} \emph{I can describe the rides, shows and stalls of a fair.}

Complete the last lesson of chapter 311: I can describe the rides, shows and stalls of a fair.
\end{chapteropening}

\section[\texorpdfstring{\bn{নাগরদোলা} (nāgôrdolā)}{নাগরদোলা (nāgôrdolā)}]{\bn{নাগরদোলা} (nāgôrdolā) — a Ferris wheel; a big wheel at a fair}

\label{lesson:BN-C311-nagordola}



\begin{quote}
Before the new one: say the Bengali for a goddess, then the Bengali for holy.
\end{quote}

\subsection*{You'll want to know: \bn{নাগরদোলা}}
\textbf{\bn{নাগরদোলা}} — \emph{nāgôrdolā} — \textquotedblleft{}a Ferris wheel; a big wheel at a fair\textquotedblright{}.

\begin{quote}
\bn{মেলায়} \bn{শিশুরা} \bn{নাগরদোলায়} \bn{চড়ল।}
\end{quote}

\begin{quote}
\emph{melāy shishurā nāgôrdolāy chôṛlô.}
\end{quote}

\begin{quote}
At the fair the children rode the big wheel.
\end{quote}

\begin{cousinweb}
\textbf{\bn{নাগরদোলা}} (\emph{nāgôrdolā}) is usually explained as \bn{নাগর}, 'of the town' — Sanskrit \emph{nāgara}, 'town-born' (Monier-Williams) — with \bn{দোলা}, a swing, the word behind \textbf{\bn{দোলানো}} (CDIAL 6453): the great swing of the fair.
\end{cousinweb}

\begin{grammarlens}[title={What you've built}]
One more word for telling what happened.
\end{grammarlens}

\subsection*{Your turn}
\begin{itemize}
\raggedright
  \item \emph{Say it:} \emph{nāgôrdolā}
  \item \emph{Say it:} \emph{nāgôrdolā}, once more
  \item \emph{Say it:} \emph{melāy shishurā nāgôrdolāy chôṛlô.} — At the fair the children rode the big wheel.
  \item \emph{Recall:} say the Bengali for a goddess, then the Bengali for holy, then say \emph{nāgôrdolā} again
\end{itemize}

\begin{tcolorbox}[breakable,colback=teal!4,colframe=teal!35!black,title={Before you move on}]
What does \bn{নাগরদোলা} mean? (\textquotedblleft{}A Ferris wheel; a big wheel at a fair\textquotedblright{}.) Say it once more.
\end{tcolorbox}
\section[\texorpdfstring{\bn{পুতুলনাচ} (putulnāch)}{পুতুলনাচ (putulnāch)}]{\bn{পুতুলনাচ} (putulnāch) — a puppet show}

\label{lesson:BN-C311-putulnach}



\begin{quote}
Before the new one: say the Bengali for holy, then the Bengali for a Ferris wheel.
\end{quote}

\subsection*{You'll want to know: \bn{পুতুলনাচ}}
\textbf{\bn{পুতুলনাচ}} — \emph{putulnāch} — \textquotedblleft{}a puppet show\textquotedblright{}.

\begin{quote}
\bn{মেলায়} \bn{সন্ধ্যায়} \bn{একটা} \bn{পুতুলনাচ} \bn{হল।}
\end{quote}

\begin{quote}
\emph{melāy sôndhyāy ekṭā putulnāch hôlô.}
\end{quote}

\begin{quote}
There was a puppet show at the fair in the evening.
\end{quote}

\begin{cousinweb}
\textbf{\bn{পুতুলনাচ}} (\emph{putulnāch}) is built from \textbf{\bn{পুতুল}}, a doll, and \textbf{\bn{নাচ}}, a dance: dolls that dance on strings.
\end{cousinweb}

\begin{grammarlens}[title={What you've built}]
One more word for telling what happened.
\end{grammarlens}

\subsection*{Your turn}
\begin{itemize}
\raggedright
  \item \emph{Say it:} \emph{putulnāch}
  \item \emph{Say it:} \emph{putulnāch}, once more
  \item \emph{Say it:} \emph{melāy sôndhyāy ekṭā putulnāch hôlô.} — There was a puppet show at the fair in the evening.
  \item \emph{Recall:} say the Bengali for holy, then the Bengali for a Ferris wheel, then say \emph{putulnāch} again
\end{itemize}

\begin{tcolorbox}[breakable,colback=teal!4,colframe=teal!35!black,title={Before you move on}]
What does \bn{পুতুলনাচ} mean? (\textquotedblleft{}A puppet show\textquotedblright{}.) Say it once more.
\end{tcolorbox}
\section[\texorpdfstring{\bn{বেলুন} (belun)}{বেলুন (belun)}]{\bn{বেলুন} (belun) — a balloon}

\label{lesson:BN-C311-belun}



\begin{quote}
Before the new one: say the Bengali for a Ferris wheel, then the Bengali for a puppet show.
\end{quote}

\subsection*{You'll want to know: \bn{বেলুন}}
\textbf{\bn{বেলুন}} — \emph{belun} — \textquotedblleft{}a balloon\textquotedblright{}.

\begin{quote}
\bn{বাবা} \bn{আমাকে} \bn{মেলা} \bn{থেকে} \bn{একটা} \bn{লাল} \bn{বেলুন} \bn{কিনে} \bn{দিলেন।}
\end{quote}

\begin{quote}
\emph{bābā āmāke melā theke ekṭā lāl belun kine dilen.}
\end{quote}

\begin{quote}
Father bought me a red balloon at the fair.
\end{quote}

\begin{cousinweb}
\textbf{\bn{বেলুন}} (\emph{belun}) is borrowed from English \emph{balloon}.
\end{cousinweb}

\begin{grammarlens}[title={What you've built}]
One more word for telling what happened.
\end{grammarlens}

\subsection*{Your turn}
\begin{itemize}
\raggedright
  \item \emph{Say it:} \emph{belun}
  \item \emph{Say it:} \emph{belun}, once more
  \item \emph{Say it:} \emph{bābā āmāke melā theke ekṭā lāl belun kine dilen.} — Father bought me a red balloon at the fair.
  \item \emph{Recall:} say the Bengali for a Ferris wheel, then the Bengali for a puppet show, then say \emph{belun} again
\end{itemize}

\begin{tcolorbox}[breakable,colback=teal!4,colframe=teal!35!black,title={Before you move on}]
What does \bn{বেলুন} mean? (\textquotedblleft{}A balloon\textquotedblright{}.) Say it once more.
\end{tcolorbox}
\section[\texorpdfstring{\bn{ঘোরা} (ghorā)}{ঘোরা (ghorā)}]{\bn{ঘোরা} (ghorā) — to go round; to wander about}

\label{lesson:BN-C311-ghora}



\begin{quote}
Before the new one: say the Bengali for a puppet show, then the Bengali for a balloon.
\end{quote}

\subsection*{You'll want to know: \bn{ঘোরা}}
\textbf{\bn{ঘোরা}} — \emph{ghorā} — \textquotedblleft{}to go round; to wander about\textquotedblright{}.

\begin{quote}
\bn{আমরা} \bn{বিকেলে} \bn{মেলায়} \bn{অনেক} \bn{ঘুরলাম।}
\end{quote}

\begin{quote}
\emph{āmrā bikele melāy ônek ghurlām.}
\end{quote}

\begin{quote}
We spent the afternoon wandering round the fair.
\end{quote}

Past (he/she …): \bn{সেই} \bn{চাকা} \bn{ঘুরল}, \emph{sei chākā ghurlô}.

The stem vowel rises in the past: \bn{ঘুরল}, turned; \bn{ঘুরলাম}, we wandered.

\begin{cousinweb}
\textbf{\bn{ঘোরা}} (\emph{ghorā}) is filed in CDIAL 4488 under a Middle Indo-Aryan \emph{ghurati}, 'revolves', which lists Bengali \emph{ghurā}, 'to revolve'. A wheel \bn{ঘোরে}; so does a visitor at a fair.
\end{cousinweb}

\begin{grammarlens}[title={What you've built}]
One more word for telling what happened.
\end{grammarlens}

\subsection*{Your turn}
\begin{itemize}
\raggedright
  \item \emph{Say it:} \emph{ghorā}
  \item \emph{Say it:} \emph{ghorā}, once more
  \item \emph{Say it:} \emph{āmrā bikele melāy ônek ghurlām.} — We spent the afternoon wandering round the fair.
  \item \emph{Recall:} say the Bengali for a puppet show, then the Bengali for a balloon, then say \emph{ghorā} again
\end{itemize}

\begin{tcolorbox}[breakable,colback=teal!4,colframe=teal!35!black,title={Before you move on}]
What does \bn{ঘোরা} mean? (\textquotedblleft{}To go round; to wander about\textquotedblright{}.) Say it once more.
\end{tcolorbox}
\section[\texorpdfstring{\bn{ঠাসা} (ṭhāsā)}{ঠাসা (ṭhāsā)}]{\bn{ঠাসা} (ṭhāsā) — packed; crammed}

\label{lesson:BN-C311-thasa}



\begin{quote}
Before the new one: say the Bengali for a balloon, then the Bengali for to go round.
\end{quote}

\subsection*{You'll want to know: \bn{ঠাসা}}
\textbf{\bn{ঠাসা}} — \emph{ṭhāsā} — \textquotedblleft{}packed; crammed\textquotedblright{}.

\begin{quote}
\bn{মেলার} \bn{দিন} \bn{বাসগুলো} \bn{লোকে} \bn{ঠাসা} \bn{ছিল।}
\end{quote}

\begin{quote}
\emph{melār din bāsgulo loke ṭhāsā chhilô.}
\end{quote}

\begin{quote}
On the day of the fair the buses were packed with people.
\end{quote}

\begin{cousinweb}
\textbf{\bn{ঠাসা}} (\emph{ṭhāsā}) is filed in CDIAL 5499 under a Middle Indo-Aryan \emph{ṭhass-}, 'press down', which lists Bengali \emph{thāsā}, 'to knead, cram'; \bn{ঠাসা} is the crammed result.
\end{cousinweb}

\begin{grammarlens}[title={What you've built}]
One more word for telling what happened.
\end{grammarlens}

\subsection*{Your turn}
\begin{itemize}
\raggedright
  \item \emph{Say it:} \emph{ṭhāsā}
  \item \emph{Say it:} \emph{ṭhāsā}, once more
  \item \emph{Say it:} \emph{melār din bāsgulo loke ṭhāsā chhilô.} — On the day of the fair the buses were packed with people.
  \item \emph{Recall:} say the Bengali for a balloon, then the Bengali for to go round, then say \emph{ṭhāsā} again
\end{itemize}

\begin{tcolorbox}[breakable,colback=teal!4,colframe=teal!35!black,title={Before you move on}]
What does \bn{ঠাসা} mean? (\textquotedblleft{}Packed; crammed\textquotedblright{}.) Say it once more.
\end{tcolorbox}
